%========== Packages ==========
\usepackage%[disable]
{todonotes}

\usepackage{graphicx}
\usepackage{braket}
\usepackage{tikz}
\usetikzlibrary{arrows, decorations.markings, decorations.pathmorphing, decorations.pathreplacing, shapes.misc, shapes.geometric, arrows.meta, fit, patterns}
\usetikzlibrary{calc}

\usepackage{amssymb}
\usepackage{stmaryrd}
\usepackage{proof}

\usepackage{mathrsfs}

\usepackage{mathtools}
\usepackage{ifthen}

\usepackage{cmll}
\usepackage{xspace}
\usepackage{adjustbox}
\adjustboxset{max width = \linewidth, center = \linewidth}
\usepackage{ebproof} % Proof trees
\ebproofset{
	template={$\vdash \inserttext$},
	right label template={$\scriptstyle(\inserttext)$},
	left label template={$\scriptstyle\inserttext$},
	center=false,
}
\usepackage{url}
%\hypersetup{colorlinks,linkcolor={red},citecolor={blue},urlcolor={red}}

%========== Work to be done ==========
\NewDocumentCommand{\ntodo}{O{}mm}{\todo[#1]{#2: #3}}
%\RenewDocumentCommand{\ntodo}{O{}m}{}
\NewDocumentCommand{\remi}{O{}m}{\ntodo[color=olive!50,#1]{R}{#2}}
\newcommand\todoc[2][]{\todo[color=pink!20,#1]{#2}} %claudia's notes
\newcommand\todocf[1]{\todoc[inline]{#1}}%claudia's notes

\newcommand{\cf}[1]{\color{violet}{#1}}
\newcommand{\CF}[1]{{\color{violet}{**CF: #1}}}

\newcommand{\pink}[1]{{\color{magenta}{#1}}}
\newcommand{\violet}[1]{{\color{violet}{#1}}}
\newcommand{\RED}[1]{{\color{red}{#1}}}
\newcommand{\blue}[1]{{\color{blue}{#1}}}

%========== CREF NAMES ==========
\Crefname{equation}{Eq.}{Eqs.}
\Crefname{section}{Sec.}{Sections}
\Crefname{figure}{Fig.}{Figs.}
\Crefname{theorem}{Thm.}{Thm.}
\Crefname{proposition}{Prop.}{Prop.}
\Crefname{remark}{Rem.}{Rem.}
\Crefname{corollary}{Cor.}{Cor.}
\Crefname{definition}{Def.}{Def.}
\Crefname{example}{Ex.}{Ex.}
\Crefname{property}{Property}{Properties}

%========== Adjustbox: centered and no larger than line width ==========
\adjustboxset{max width = \linewidth, center = \linewidth}

%========== Tikz ==========
% General scope
% 1st optional argument: color
% 2nd optional argument: edge style
% 3rd optional argument: line width (thick = 0.8pt)
\NewDocumentEnvironment{genscope}{O{black}O{solid,->}O{0.8pt}}{
\begin{scope}[
every node/.style={circle,minimum size=8mm,inner sep=0,draw=#1,line width=#3},
every edge/.style={draw=#1,#2,line width=#3},
every path/.style={thick,draw=#1,#2,-,line width=#3}]
\color{#1}
}{
\end{scope}
}

% Style for naming edges and paths
\tikzset{
    named edge/.style={draw=none,inner sep=1pt,outer sep=2pt,rectangle,minimum size=0pt,auto},
    named path/.style={named edge,sloped},
}

%========== Abbreviations ==========
\newcommand*{\resp}{resp.}
\newcommand*{\ie}{\emph{i.e.}\xspace}
%\newcommand*{\wolog}{\emph{w.l.o.g.}}
\newcommand*{\eg}{\emph{e.g.}\xspace}
%\newcommand*{\cf}{\emph{c.f.}}
\newcommand*{\aka}{\emph{a.k.a.}}
\newcommand*{\wrt}{\emph{w.r.t.}}

%========== Notations ==========
% in-text definition
\newcommand{\textdef}[1]{\textbf{#1}}
% Emphase when defining a term
\newcommand*{\definitive}[1]{\textbf{#1}}

\newcommand*{\defeq}{\stackrel{\text{def}}{=}}

%========== General ==========
% cardinality of a set
\newcommand{\Card}[1]{\#(#1)}
% Typing
\newcommand{\der}{\vdash}
% Symmetric difference
\newcommand*{\symdif}{\mathrel{\Delta}}
% Difference of sets
\newcommand*{\difset}{\mathrel{-}}

\newcommand{\true}{\mathtt{t}}
\newcommand{\false}{\mathtt{f}}
\newcommand{\Real}{\mathbb{R}}
\newcommand{\bool}{\mathbb{B}}

%========== Linear Logic ==========
% Dual
\newcommand*{\orth}{^\perp}
% Tensor
\newcommand*{\tens}{\otimes}
% Linear Implication
\newcommand*{\lollipop}{\multimap}
% 1
\newcommand*{\one}{1}
% 0
%\newcommand*{\zero}{0}

\newcommand{\Vars}{\mathsf{Names}} %set of propositional variables
\newcommand{\VarA}{X} %metavariable of a propositional variable 
\newcommand{\VarB}{Y} %metavariable of a propositional variable 
\newcommand{\VarC}{Z} %metavariable of a propositional variable 
\newcommand{\One}{\mathsf 1} %unit tensor
\newcommand{\Bottom}{\mathsf \perp} %unit par

\newcommand{\FormA}{F} %metavariable of MLL formula 1
\newcommand{\FormB}{G} %metavariable of MLL formula 2
\newcommand{\FormC}{H} %metavariable of MLL formula 3

% MLL fragment of linear logic
\newcommand{\MLL}{\ensuremath{\mathsf{MLL}}\xspace}

% Positive and Negative atoms
\newcommand*{\posatom}[1]{{#1}^+}
\newcommand*{\negatom}[1]{{#1}^-}

% Name of rules with text
\newcommand*{\rulenametext}[1]{\textnormal{\textit{#1}}}
\newcommand*{\ax}{\rulenametext{ax}}
\newcommand*{\cut}{\rulenametext{cut}}
\newcommand*{\mix}{\rulenametext{mix}}
\newcommand*{\ex}{\rulenametext{ex}}

\newcommand{\Xn}{X^-}
\newcommand{\Yn}{Y^-}
\newcommand{\Zn}{Z^-}
\newcommand{\X}{X^+}
\newcommand{\Y}{Y^+}
\newcommand{\Z}{Z^+}

% Names in a formula/sequent
\newcommand{\Nm}[1]{\mathsf{Nm}(#1)}

%========== Proof-Nets ==========
% Box of a vertex
\newcommand*{\boxof}[1]{\ensuremath{\bbox^{#1}}}
% Proof-net of a Bayesian network
\newcommand*{\pnof}[1]{\mathcal{P}(#1)}

\newcommand*{\logsys}[1]{\textnormal{\textsf{#1}}}
% MLL fragment
\renewcommand*{\MLL}{\ensuremath{\logsys{MLL}}\xspace}
% MLL with boxes
\newcommand*{\MLLB}{\ensuremath{\logsys{MLL}^{\mathtt{B}}}}

% Some proof-net
\newcommand*{\somepn}{\mathcal{R}}
% Some node
\newcommand*{\somenode}{n}
% Some graph
\newcommand*{\G}{\mathcal{G}}

\newcommand{\CutS}[1]{\mathrm{Cut}(#1)}
\newcommand{\Cuts}[1]{\mathrm{Cut}(#1)}

\newcommand{\netN}{\mathcal N}
\newcommand{\N}{\netN}
\newcommand{\Net}{\N}
\newcommand{\netR}{\mathcal R}
\newcommand{\R}{\netR}
\newcommand{\ProofN}{\R}
\newcommand{\netM}{\mathcal M}
\newcommand{\netG}{\mathcal G}
\newcommand{\netC}{\mathcal C}
\newcommand{\netS}{\mathcal S}
\newcommand{\netT}{\mathcal T}
\newcommand{\netE}{\mathcal E}
\newcommand{\E}{\netE}

\newcommand{\Module}{\M}

\newcommand{\contr}{\cn}

\newcommand{\bbox}{\mathtt{b}}
\newcommand{\NodeB}{\bbox}

\newcommand{\Boxs}[1]{\mathsf{Box}(#1)}
\newcommand{\Boxes}[1]{\mathsf{Boxes}(#1)}
\newcommand{\boxn}{\mathsf{box}}
\newcommand{\weak}{\mathsf{w}}
\newcommand{\cn}{\mathtt{c}}

\newcommand{\Out}[2]{\mathsf{Split}_{#2}({#1})}
\newcommand{\Pull}[2]{\mathsf{Pull}_#2(#1)}
\newcommand*{\pullingout}{isolation}
\newcommand*{\pulledout}[2]{\ensuremath{\Out{#1}{#2}}}

% Polarized part of a proof-net
\newcommand{\pol}[1]{\mathit{Pol}({#1})}

%========== Abstract Rewriting System ==========
\newcommand{\rules}{\mathtt{R}}
\newcommand{\rr}{\mathtt{r}}
% Reduction step to
\newcommand{\Red}[1][]{\xrightarrow{#1}}
% Reduction step from
\newcommand{\FromRed}[1][]{\xleftarrow{#1}}
\newcommand{\Reds}[1][]{\mathrel{\mathop{\xrightarrow{#1}}\!{}^\ast}} %reflexive-transitive closure of a reduction step
\newcommand{\FromReds}[1][]{\mathrel{\mathop{{}^\ast\xleftarrow{#1}}}} %reflexive-transitive closure of a reduction step from
\newcommand{\Redsplus}[1][]{\mathrel{\mathop{\xrightarrow{#1}}\!{}^+}} %transitive closure of a reduction step

%========== Semantics ==========
\newcommand{\den}[1]{\llbracket{#1}\rrbracket}
\newcommand{\sem}[1]{\llparenthesis{#1}\rrparenthesis}
% Operations in the semantics
\newcommand*{\compfact}{\odot}

% Category Q of Selinger
\newcommand*{\Q}{\ensuremath{\mathbf{Q}}}

%========== Classical ==========
% Type of booleans
\newcommand*{\booltype}{\mathbb{B}}
% Non-negative reals
\newcommand{\nnReal}{\Real^+_0}
% Set of a random variable
\newcommand*{\setofrv}[1]{\mathsf{Val}(#1)}
% Probability of
\newcommand*{\pb}[1]{Pr(#1)}
% Set of values of a random variable
\newcommand{\Val}[1]{\mathsf{Val}(#1)}

% Unit factor
\newcommand{\ftone}{\mathbf{1}}

\newcommand{\Fprod}{\odot} %product of factor (corresponding to tensor product + some diagonalisations), infix notation
\newcommand{\BigFProd}{\bigodot} %n-ary product of factor (corresponding to tensor product + some diagonalisations)
\newcommand{\FSum}[2]{\textstyle\sum_{#1}(#2)} %summing out a rv from a factor (corresponding to apply weakening), prefix notation
\newcommand{\FProj}[2]{\mathrm{proj}_{#2}\left(#1\right)} %projection of a factor #1 to a set of rvs #2

% Some Bayesian network
\newcommand*{\somebayes}{\mathcal{B}}
\newcommand*{\bn}{\somebayes}

\newcommand{\bX}{\mathbf X}
\newcommand{\bY}{\mathbf Y}
\newcommand{\bZ}{\mathbf Z}
\newcommand{\bS}{\mathbf S}
\newcommand{\bE}{\mathbf E}
\newcommand{\bC}{\mathbf C}
\newcommand{\bW}{\mathbf W}
\newcommand{\bV}{\mathbf V}

\newcommand{\bx}{{\mathbf x}}
\newcommand{\by}{{\mathbf y}}
\newcommand{\bz}{{\mathbf z}}
\newcommand{\ba}{{\mathbf a}}
\newcommand{\bb}{{\mathbf b}}
\newcommand{\bc}{{\mathbf c}}
\newcommand{\bd}{{\mathbf d}}
\newcommand{\be}{{\mathbf e}}

%========== Quantum ==========
% Trace
\newcommand{\tr}[1]{\mathrm{tr}\left(#1\right)}
% Partial trace
\newcommand{\ptr}[1]{\mathrm{tr}_{#1}}

% Conjugate-transpose of a matrix
\newcommand*{\ct}{^\ast}

% Set of linear operators
\renewcommand{\L}{\mathcal{L}}
% Set of positive operators
\newcommand{\Pos}[1]{{\L(#1)}^+}
% Set of density matrices
\newcommand{\D}{\mathcal{D}}
\newcommand{\DM}{\D}

\newcommand{\Eop}{\mathcal E}

\newcommand{\U}{\mathcal{U}}

% Canonical basis of a hilbert space C^{n x n}
\newcommand*{\baseof}[1]{\textsf{Basis}_{#1}}

% Complex
\newcommand*{\C}{\mathbb{C}}

% Hilbert space
\newcommand{\HH}{\mathcal{H}}

% Functional view of a quantum factor
\newcommand{\qfactfunc}[1]{\widehat{#1}}

% Some density matrix
\newcommand{\M}{\rho}
\newcommand{\dm}{\rho}

% Some quantum register
\newcommand{\reg}{Q}
\newcommand{\qreg}{\reg}

% Some set of quantum registers
\newcommand{\bQ}{\mathbf Q}
% Some quantum factor
\newcommand*{\someqfact}{\phi}
\newcommand*{\someqfactset}{\Phi}

%% Linear functions
% id
\newcommand*{\id}[1][]{\mathrm{Id}_{#1}}
% Cup
\newcommand*{\qcup}[1][]{\mathrm{cup}_{#1}}
% Cap
\newcommand*{\qcap}[1][]{\mathrm{cap}_{#1}}
% Contraction
\newcommand*{\cc}[1][]{\mathrm{c}_{#1}}
% Weakening
\newcommand*{\cw}[1][]{\mathrm{w}_{#1}}

% The identity on density matrices
\newcommand{\idop}{\id}

% Choi-Jamiolkowski isomorphism
\newcommand*{\choiiso}{\mathcal{J}}

%========== Nomenclature ==========
\newcommand{\rv}{r.v.\xspace}
\newcommand{\rvs}{r.v.s\xspace}

% CPTs
\newcommand{\CPT}{CPT\xspace}

\newcommand{\dags}{dags}

\newcommand{\BN}{Bayesian network\xspace}
\newcommand{\BNs}{Bayesian networks\xspace}

% Quantum factors & Q-CPTs
\newcommand*{\qfactor}{Q-factor}
\newcommand*{\Qfactor}{\qfactor}

\newcommand*{\QCPT}{Quantum CPT}

\newcommand*{\qCPT}{Q-\textsf{cpt}\xspace}
\newcommand*{\qCPTs}{Q-\textsf{cpt}s\xspace}

% Quantum Bayesian Networks
\newcommand*{\QBN}{quantum Bayesian network}
\newcommand*{\QBNs}{Quantum Bayesian networks\xspace}

\newcommand*{\qqbn}{\textsf{QBN}}

% Quantum (Bayesian) Proof-nets
\newcommand*{\qpn}{quantum proof-net\xspace}
\newcommand*{\qpns}{quantum proof-nets\xspace}
\newcommand*{\Qpn}{Quantum proof-net\xspace}
\newcommand*{\Qpns}{Quantum proof-nets\xspace}

\newcommand*{\qbpn}{\textsf{qpn}\xspace}
\newcommand*{\qbpns}{\textsf{qpns}\xspace}

\newcommand*{\main}{output\xspace}

\newcommand*{\good}{closed\xspace}

\newcommand{\cpts}{CPT's\xspace}
\newcommand{\cpt}{CPT\xspace}

\newcommand{\pn}{proof-net\xspace}
\newcommand{\pns}{proof-nets\xspace}

\newcommand{\pt}{proof-tree\xspace}
\newcommand{\pts}{proof-trees\xspace}
